\documentclass[11pt]{article}
% Publication-quality resume template (LuaLaTeX + latexmk recommended)
% US Letter, single-column, restrained typography

\usepackage{ifxetex}
\usepackage{ifluatex}
\ifnum\ifxetex 1\fi\ifluatex 1\fi=0
  \usepackage[T1]{fontenc}
  \usepackage[utf8]{inputenc}
\else
  \usepackage{fontspec}
  \defaultfontfeatures{Ligatures=TeX,Scale=MatchLowercase}
  % Reproducible fonts available in TeX Live
  \setmainfont{TeX Gyre Termes}
  \setsansfont{TeX Gyre Heros}
  \setmonofont{TeX Gyre Cursor}
\fi

\usepackage{microtype}
\usepackage[letterpaper,margin=0.7in]{geometry}
\usepackage{parskip}
\usepackage[hidelinks]{hyperref}
\usepackage{xcolor}
\usepackage{enumitem}
\setlist[itemize]{noitemsep, topsep=2pt, parsep=0pt, itemsep=2pt, leftmargin=*, label=\textbullet}
\usepackage{ragged2e}
% Optional TikZ emblems (kept minimal; not used for core text)
\usepackage{tikz}

% Centralized typographic sizes
\newcommand{\namefontsize}{25pt}
\newcommand{\bodyfontsize}{11pt}
\newcommand{\sectionfontsize}{12pt}
\newcommand{\contactfontsize}{10pt}
\renewcommand{\baselinestretch}{1.18} % roughly 13pt leading for 11pt body
\makeatletter
\ifx\@setfontsize\undefined\else\fi
\makeatother

% Pandoc macro definitions
\providecommand{\tightlist}{%\setlength{\itemsep}{0pt}\setlength{\parskip}{0pt}}%

% Header layout commands
\newcommand{\ResumeName}[1]{\vspace*{-0.5em}\noindent{\sffamily\bfseries\fontsize{\namefontsize}{0}\selectfont #1}\vspace{0.2em}\\}
\newcommand{\ContactLine}[1]{\noindent{\small\normalfont #1}\vspace{0.1em}\\}

% Section heading
\usepackage{sectsty}
\sectionfont{\sffamily\bfseries\fontsize{\sectionfontsize}{0}\selectfont}
\subsectionfont{\sffamily\bfseries\fontsize{11pt}{0}\selectfont}

% Ensure hyphenation and ragged-right for better extraction
\raggedright

\begin{document}
\pagestyle{plain}

% Title block from metadata: use Pandoc's Marco Del Fava, Marco Del
Fava, and custom true metadata
  \ResumeName{Marco Del Fava}
      \noindent\small\texttt{mdelfava@ucdavis.edu} \quad 707-637-6544 \\
    \href{https://github.com/BarkBarkBarkBarkBarkBarkBark}{GitHub}\quad 
    \href{www.linkedin.com/in/marc-del-fava}{LinkedIn}\quad 
  
% Main body from pandoc
\subsection{Summary}\label{summary}

Recently admitted to the UC Davis M.S. in Biomedical Engineering (Fall
2026) and active neuroscience researcher working at the intersection of
human intracranial electrophysiology, single-unit digital signal
processing, and neurotechnology tooling. Experience includes spike
sorting with Kilosort4, O-Sort, and SpikeInterface; building
GPU-accelerated pipelines on AWS; and implementing Dark Horse Neuro /
Neuralynx acquisition workflows for replay, compression-aware MED
recording, and large-scale neural data handling. I look forward to a
career in neural signal processing, neural acquisition systems, and
BCI-related infrastructure.

\subsection{Selected Projects}\label{selected-projects}

\subsubsection{Dark Horse Neuro / Neuralynx acquisition pipeline ---
2026}\label{dark-horse-neuro-neuralynx-acquisition-pipeline-2026}

Links:
\href{https://github.com/BarkBarkBarkBarkBarkBarkBark/dhn_client}{GitHub}
·
\href{https://barkbarkbarkbarkbarkbarkbark.github.io/dhn_client/}{Docs}
- Built a Python orchestration layer to launch DHN-AQ headlessly over
SSH, generate RC/CS session files, and coordinate Pegasus / NetCom
replay workflows for human intracranial recordings. - Implemented an NRD
replay bridge that re-broadcasts ATLAS-format UDP packets so archived
Neuralynx sessions can be ingested by DHN-AQ as if hardware were live. -
Automated compression-aware MED recording setup for 512-channel macro /
micro configurations, including decimation, lossless/lossy options, and
pre-flight network diagnostics.

\subsubsection{spike\_discrim single-unit DSP workbench ---
2026}\label{spike_discrim-single-unit-dsp-workbench-2026}

Links:
\href{https://github.com/BarkBarkBarkBarkBarkBarkBark/spike_discrim}{GitHub}
- Benchmarked low-latency waveform representations for real-time
single-unit discrimination using curated OSort snippets, synthetic data,
AUC, balanced accuracy, mutual information, and KNN purity. - Built a
Numba-accelerated WeightBank population code to score spike-likeness and
project waveforms into afferent activation space for clustering and
unit-separation analysis.

\subsubsection{SNN Agent real-time spike sorting ---
2026}\label{snn-agent-real-time-spike-sorting-2026}

Links:
\href{https://github.com/BarkBarkBarkBarkBarkBarkBark/SNNeuro}{GitHub} -
Developed a streaming spiking neural network with bandpass filtering,
Kalman noise gating, competitive STDP, and browser-based monitoring for
live spike detection and sorting. - Evaluated 20 kHz real-time pipeline
behavior across synthetic, .ncs, UDP, and LSL inputs with online
decoding for closed-loop neurotech use cases.

\subsection{Experience}\label{experience}

\textbf{Operations Specialist} --- \emph{First Step Communities} ---
Sacramento, CA Jul 2021 --- Present

\begin{itemize}
\tightlist
\item
  Analyzed program and partner datasets using Python, SPSS, and Excel to
  identify predictors of successful exits and service effectiveness.
\item
  Coordinated data collection from clinical partners including UC Davis
  Health and Kaiser Permanente to support integrated care outcome
  analyses.
\item
  Identified an association between on-site clinical engagement and
  permanent housing outcomes; supported leadership in scaling integrated
  clinical services across sites.
\item
  Developed 40+ standardized data collection forms and automated
  reporting workflows using Microsoft Power Automate.
\end{itemize}

\textbf{Analyst} --- \emph{Eurofins Test America} --- West Sacramento,
CA Oct 2019 --- Jun 2021

\begin{itemize}
\tightlist
\item
  Maintained high-accuracy data entry and review in a regulated
  laboratory environment; supported QA documentation and adherence to
  standard operating procedures.
\item
  Performed solvent extraction of organic compounds from soil and water
  samples in a high-throughput production setting.
\item
  Coordinated daily scheduling and backlog management to meet turnaround
  targets across teams.
\end{itemize}

\subsection{Education}\label{education}

\textbf{M.S. Biomedical Engineering}, University of California, Davis,
Davis, CA --- Beginning Fall 2026 Recently admitted; incoming first-year
master's student

\textbf{Python for Data Science, Professional Concentration}, UC Davis
Continuing and Professional Education, Davis, CA --- 2025

\textbf{M.S. Psychology}, Arizona State University, Tempe, AZ --- 2023

\textbf{B.S. Neuroscience, Physiology, and Behavior}, University of
California, Davis, Davis, CA --- 2017

\subsection{Skills}\label{skills}

\textbf{Neural Data \& Signal Processing:} Single-unit DSP and waveform
feature extraction, Spike sorting (Kilosort4, O-Sort, SpikeInterface),
Neural data processing (filtering, PCA, template matching, clustering),
Intracranial electrophysiology (microwire, LFP, iEEG)

\textbf{Neural Acquisition \& Tooling:} Neuralynx Pegasus / ATLAS /
NetCom workflows, Dark Horse Neuro DHN-AQ RC/CS configuration, NRD
replay, UDP relay, and MED archive generation, Linux, SSH automation,
xvfb, and acquisition diagnostics

\textbf{Machine Learning \& Data Science:} Python (NumPy, Pandas,
scikit-learn, PyTorch, Numba), Statistical analysis (SPSS), SQL and
database engineering (PostgreSQL, PGVector), Cloud GPU computing (AWS
EC2, CUDA, Docker) and MATLAB workflows

\textbf{General Technical:} Git and version control, Linux and shell
scripting, Web interfaces (Streamlit), Microsoft 365 administration

\textbf{Data Systems \& Reporting:} Dashboards and stakeholder reporting
(Looker), Workflow automation (Microsoft 365, Power Automate), Survey
design and focus group facilitation, Equity-informed outcomes
measurement

\subsection{Certifications}\label{certifications}

\begin{itemize}
\tightlist
\item
  \textbf{Data Analysis with Matplotlib, Pandas, and Scikit-Learn Python
  Libraries} --- UC Davis Continuing and Professional Education
  (2025-09-18) ---
  \href{https://www.credly.com/badges/505ef11c-58b4-48bd-aa9e-5a21f0247d46/public_url}{Badge}-
  \textbf{CITI --- Biomedical Research} --- CITI Program- \textbf{CITI
  --- Essentials of Research Administration} --- CITI Program-
  \textbf{CITI --- GCP Social and Behavioral Research Best Practices}
  --- CITI Program Recently admitted to the UC Davis M.S. in Biomedical
  Engineering (Fall 2026) and active neuroscience researcher working at
  the intersection of human intracranial electrophysiology, single-unit
  digital signal processing, and neurotechnology tooling. Experience
  includes spike sorting with Kilosort4, O-Sort, and SpikeInterface;
  building GPU-accelerated pipelines on AWS; and implementing Dark Horse
  Neuro / Neuralynx acquisition workflows for replay, compression-aware
  MED recording, and large-scale neural data handling. I look forward to
  a career in neural signal processing, neural acquisition systems, and
  BCI-related infrastructure.
\end{itemize}

\begin{center}\rule{0.5\linewidth}{0.5pt}\end{center}

\subsection{Neuroscience Projects}\label{neuroscience-projects}

\textbf{Dark Horse Neuro / Neuralynx acquisition pipeline} \textbar{}
2026

\begin{itemize}
\item
  Built a Python orchestration layer to launch DHN-AQ headlessly over
  SSH, generate RC/CS session files, and coordinate Pegasus / NetCom
  replay workflows for human intracranial recordings.
\item
  Implemented an NRD replay bridge that re-broadcasts ATLAS-format UDP
  packets so archived Neuralynx sessions can be ingested by DHN-AQ as if
  hardware were live.
\item
  Automated compression-aware MED recording setup for 512-channel macro
  / micro configurations, including decimation, lossless/lossy options,
  and pre-flight network diagnostics.
\item
  Links:
  \href{https://github.com/BarkBarkBarkBarkBarkBarkBark/dhn_client}{GitHub}
  ·
  \href{https://barkbarkbarkbarkbarkbarkbark.github.io/dhn_client/}{Docs}
  \textbf{spike\_discrim single-unit DSP workbench} \textbar{} 2026
\item
  Benchmarked low-latency waveform representations for real-time
  single-unit discrimination using curated OSort snippets, synthetic
  data, AUC, balanced accuracy, mutual information, and KNN purity.
\item
  Built a Numba-accelerated WeightBank population code to score
  spike-likeness and project waveforms into afferent activation space
  for clustering and unit-separation analysis.
\item
  Links:
  \href{https://github.com/BarkBarkBarkBarkBarkBarkBark/spike_discrim}{GitHub}
  \textbf{SNN Agent real-time spike sorting} \textbar{} 2026
\item
  Developed a streaming spiking neural network with bandpass filtering,
  Kalman noise gating, competitive STDP, and browser-based monitoring
  for live spike detection and sorting.
\item
  Evaluated 20 kHz real-time pipeline behavior across synthetic, .ncs,
  UDP, and LSL inputs with online decoding for closed-loop neurotech use
  cases.
\item ~
  \subsection{\texorpdfstring{Links:
  \href{https://github.com/BarkBarkBarkBarkBarkBarkBark/SNNeuro}{GitHub}}{Links: GitHub}}\label{links-github}
\end{itemize}

\subsection{Experience}\label{experience-1}

\textbf{Operations Specialist --- First Step Communities} \textbar{}
Sacramento, CA \textbar{} Jul 2021 -- Present

\begin{itemize}
\tightlist
\item
  Analyzed program and partner datasets using Python, SPSS, and Excel to
  identify predictors of successful exits and service effectiveness.
\item
  Coordinated data collection from clinical partners including UC Davis
  Health and Kaiser Permanente to support integrated care outcome
  analyses.
\item
  Identified an association between on-site clinical engagement and
  permanent housing outcomes; supported leadership in scaling integrated
  clinical services across sites.
\item
  Developed 40+ standardized data collection forms and automated
  reporting workflows using Microsoft Power Automate.
\end{itemize}

\textbf{Analyst --- Eurofins Test America} \textbar{} West Sacramento,
CA \textbar{} Oct 2019 -- Jun 2021

\begin{itemize}
\tightlist
\item
  Maintained high-accuracy data entry and review in a regulated
  laboratory environment; supported QA documentation and adherence to
  standard operating procedures.
\item
  Performed solvent extraction of organic compounds from soil and water
  samples in a high-throughput production setting.
\item
  Coordinated daily scheduling and backlog management to meet turnaround
  targets across teams.
\end{itemize}

\begin{center}\rule{0.5\linewidth}{0.5pt}\end{center}

\subsection{Research \& Volunteer
Experience}\label{research-volunteer-experience}

\textbf{Research Volunteer --- UC Davis Health --- Z-Lab (Dr.~Jie
Zheng)} \textbar{} Sacramento, CA \textbar{} Jul 2025 -- Present

\begin{itemize}
\tightlist
\item
  Build a GPU-accelerated spike sorting pipeline for dense
  microelectrode arrays using Kilosort4 on AWS EC2 instances to support
  reproducible, high-throughput reprocessing of human intracranial
  recordings.
\item
  Sort single-unit activity from human Phase II epilepsy microwire
  recordings using MATLAB (O-Sort) and Python (SpikeInterface), with an
  emphasis on preserving unit-specific waveform structure for downstream
  analysis.
\item
  Curated a dataset of approximately 2 million hand-labeled waveforms
  for spike-vs-noise modeling and benchmarking.
\item
  Trained neural network models to classify putative single units versus
  noise, improving manual curation efficiency for human recordings.
\end{itemize}

\textbf{Research Volunteer --- UC Davis Health --- Neuroscience Unit E5}
\textbar{} Sacramento, CA \textbar{} Dec 2022 -- Jul 2023

\begin{itemize}
\tightlist
\item
  Prepared headboxes and EEG application kits for the Epilepsy
  Monitoring Unit.
\item
  Shadowed EEG technologists and observed clinical neurophysiology
  procedures for patients with epilepsy and traumatic brain injury.
\end{itemize}

\textbf{Research Intern --- California Home Builders Association}
\textbar{} Sacramento, CA \textbar{} Jul 2023 -- Feb 2024

\begin{itemize}
\tightlist
\item
  Designed a cost-estimation model for emergency shelter villages with
  on-site medical facilities to inform California housing policy.
\item
  Synthesized findings into policy recommendations to support emergency
  shelter development legislation.
\end{itemize}

\subsection{Education}\label{education-1}

\textbf{M.S. Biomedical Engineering} \textbar{} University of
California, Davis \textbar{} Beginning Fall 2026 Recently admitted;
incoming first-year master's student

\textbf{Python for Data Science, Professional Concentration (GPA: 4.0)}
\textbar{} UC Davis Continuing and Professional Education \textbar{}
2025

\textbf{M.S. Psychology (GPA: 3.7)} \textbar{} Arizona State University
\textbar{} 2023

\textbf{B.S. Neuroscience, Physiology, and Behavior (GPA: 3.1)}
\textbar{} University of California, Davis \textbar{} 2017

\begin{center}\rule{0.5\linewidth}{0.5pt}\end{center}

\subsection{Relevant Coursework}\label{relevant-coursework}

\begin{itemize}
\tightlist
\item
  Neurobiology of Behavior --- UC Davis
\item
  Physiology of Psychology --- Arizona State University
\item
  Quantitative Analysis --- Arizona State University
\item
  Research Methods --- Arizona State University
\end{itemize}

\begin{center}\rule{0.5\linewidth}{0.5pt}\end{center}

\subsection{Skills}\label{skills-1}

\textbf{Neural Data \& Signal Processing:} Single-unit DSP and waveform
feature extraction, Spike sorting (Kilosort4, O-Sort, SpikeInterface),
Neural data processing (filtering, PCA, template matching, clustering),
Intracranial electrophysiology (microwire, LFP, iEEG)

\textbf{Neural Acquisition \& Tooling:} Neuralynx Pegasus / ATLAS /
NetCom workflows, Dark Horse Neuro DHN-AQ RC/CS configuration, NRD
replay, UDP relay, and MED archive generation, Linux, SSH automation,
xvfb, and acquisition diagnostics

\textbf{Machine Learning \& Data Science:} Python (NumPy, Pandas,
scikit-learn, PyTorch, Numba), Statistical analysis (SPSS), SQL and
database engineering (PostgreSQL, PGVector), Cloud GPU computing (AWS
EC2, CUDA, Docker) and MATLAB workflows

\textbf{General Technical:} Git and version control, Linux and shell
scripting, Web interfaces (Streamlit), Microsoft 365 administration

\textbf{Data Systems \& Reporting:} Dashboards and stakeholder reporting
(Looker), Workflow automation (Microsoft 365, Power Automate), Survey
design and focus group facilitation, Equity-informed outcomes
measurement

\begin{center}\rule{0.5\linewidth}{0.5pt}\end{center}

\subsection{Certifications}\label{certifications-1}

\begin{itemize}
\tightlist
\item
  \textbf{Data Analysis with Matplotlib, Pandas, and Scikit-Learn Python
  Libraries} --- UC Davis Continuing and Professional Education
  (2025-09-18) ---
  \href{https://www.credly.com/badges/505ef11c-58b4-48bd-aa9e-5a21f0247d46/public_url}{View
  badge}
\item
  CITI --- Biomedical Research
\item
  CITI --- Essentials of Research Administration
\item
  CITI --- GCP Social and Behavioral Research Best Practices
\end{itemize}

\end{document}
